\section{Proofs and the Identification Problem}
\label{app:theory}
\label{sec:appendix}

\subsection{Supporting finite-state statements}

This appendix supports the prediction diagnostics in
\cref{sec:setup}, the distinction between linear and exact
orthogonal realizations in \cref{sec:boundaries}, and the
interpretation of the experiments in \cref{sec:evidence}.
We distinguish exact realizability on a finite bank,
prediction risk in specified analytical constructions,
and error identities for fitted operators.
These supporting statements are not presented as new
general representation theory.

Let a finite group $G$ act on a finite state set $S$ through
the prescribed action $\alpha:G\to\mathrm{Sym}(S)$.
Write $g(i)=\alpha(g)(i)$.
For rows indexed by state and template, define
\[
 (P_gZ)_{(i,t)}=z_{g(i),t}.
\]
Thus
\[
 P_gP_h=P_{h\circ g}.
\]
Matrices act on representation rows by right multiplication,
so an exact row action satisfies
\[
 R_gR_h=R_{h\circ g},
 \qquad R_e=\Id.
\]
The same notation applies to a single template.
The symbol $M^\dagger$ denotes the Moore--Penrose
pseudoinverse of a matrix $M$.

\begin{proposition}[Finite-orbit realizability]
\label{prop:finite}
Let $n\geq2$ and let
$X\in\R^{n\times(n-1)}$ have zero mean across rows and
rank $n-1$.
Every permutation of its rows has a unique invertible
linear realization.
In particular, for the prescribed state action,
\[
 R_g=X^\dagger P_gX
\]
is the unique matrix satisfying $XR_g=P_gX$, and
\[
 R_gR_h=R_{h\circ g}.
\]
\end{proposition}

The proof and its consequence for second-moment whitening
are given in \cref{app:finite-orbit-proof}.

For arbitrary fitted matrices, write
\[
 e_g(i,t)=z_{i,t}R_g-z_{g(i),t}.
\]

\begin{proposition}[Prediction-error propagation]
\label{prop:propagation}
For $p=g_1\circ g_2$, the composition residual satisfies
\begin{align}
 z_{i,t}(R_{g_2}R_{g_1}-R_p)
 &=e_{g_2}(i,t)R_{g_1} \nonumber\\
 &\quad+e_{g_1}(g_2(i),t)-e_p(i,t).
\end{align}
Consequently,
\begin{align*}
 &\norm{z_{i,t}(R_{g_2}R_{g_1}-R_p)}_2\\
 &\quad\leq
 \norm{e_{g_2}(i,t)}_2\norm{R_{g_1}}_{\rm op}
 +\norm{e_{g_1}(g_2(i),t)}_2
 +\norm{e_p(i,t)}_2.
\end{align*}
\end{proposition}

The proof and the conditions needed to relate this identity
to normalized diagnostics are given in
\cref{app:propagation-proof}.

\begin{proposition}[Global-field cancellation]
\label{prop:additive}
For any permutation $g$ of a complete, uniformly weighted
state set in a template,
\[
 \sum_{i\in S}(z_{g(i),t}-z_{i,t})=0.
\]
Therefore the balanced training estimator
\[
 v_g=\mathbb E_{i,t}
       [z_{g(i),t}-z_{i,t}]
\]
is zero for every group element.
\end{proposition}

The proof and its scope are given in
\cref{app:global-cancellation-proof}.


\subsection{Finite-orbit realizability}
\label{app:finite-orbit-proof}

\begin{proof}[Proof of \cref{prop:finite}]
Centering implies
$\operatorname{col}(X)\subseteq\mathbf1^\perp$.
Both spaces have dimension $n-1$, so they are equal.
Every permutation matrix preserves $\mathbf1^\perp$;
hence $P_gX$ lies in the column space of $X$.
Full column rank gives the unique solution
\[
 R_g=(X^\top X)^{-1}X^\top P_gX.
\]
Using $P_gP_h=P_{h\circ g}$,
\[
 XR_gR_h
 =P_gXR_h
 =P_gP_hX
 =P_{h\circ g}X
 =XR_{h\circ g}.
\]
Full column rank establishes the multiplication rule.
The identity permutation gives $R_e=\Id$.
Taking $h=g^{-1}$ gives both inverse identities, so every
$R_g$ is invertible.
The argument applies to every permutation of the $n$ rows.
\end{proof}

If $n$ lexical vectors are affinely independent, their
centered span has dimension $n-1$.
Retaining all these PCA dimensions gives coordinates on
that entire span, so the proposition applies.
Neither regular polygon geometry nor prior knowledge of a
named group is required.
The resulting matrices need not be orthogonal in the
original Euclidean coordinates.

For the uniform state distribution, however, define
\[
 \Sigma=\frac1nX^\top X\succ0.
\]
Exact realization and $P_g^\top P_g=\Id_n$ imply
\[
 R_g^\top\Sigma R_g
 =\frac1n(P_gX)^\top(P_gX)
 =\Sigma.
\]
In the whitened coordinates
\[
 \bar X=X\Sigma^{-1/2},
 \qquad
 Q_g=\Sigma^{1/2}R_g\Sigma^{-1/2},
\]
we therefore have
\[
 \bar XQ_g=P_g\bar X,
 \qquad
 Q_g^\top Q_g=\Id_{n-1},
 \qquad
 Q_gQ_h=Q_{h\circ g}.
\]
Thus every state permutation has an exact orthogonal
realization after whitening this full centered state span.
This is the shared-code case used in
\cref{sec:boundaries}; it does not assert that arbitrary
truncated or context-dependent coordinates have the same
property.

A related interpolation issue arises in representation
comparison: invariance to arbitrary invertible linear
changes can make sufficiently high-dimensional sample
representations indistinguishable
\citep{kornblith2019similarity}.
That result concerns comparison statistics in the relevant
sample-rank regime.
It does not establish transformation prediction in new
contexts or imply that the present 16-dimensional banks
fall into that regime.


\subsection{Separate contextual realizations need not be shared}

Consider the centered three-state coordinates
\[
 X=
 \begin{bmatrix}
  1&0\\
  0&1\\
  -1&-1
 \end{bmatrix},
 \qquad
 R=
 \begin{bmatrix}
  0&1\\
  -1&-1
 \end{bmatrix}.
\]
For the cycle $g:1\mapsto2\mapsto3\mapsto1$,
$XR=P_gX$.
A second contextual copy $Y=XA$, with
$A=\operatorname{diag}(2,1)$, has the forced operator
\[
 R'=A^{-1}RA
 =
 \begin{bmatrix}
  0&1/2\\
  -2&-1
 \end{bmatrix}
 \ne R.
\]
A shared operator $M$ would satisfy both
$XM=P_gX$ and $YM=P_gY$.
Full column rank forces $M=R$ from the first equation and
$M=R'$ from the second, a contradiction.

This example shows that separate realizability within
each context does not guarantee one realization in the
common coordinates.
The contradiction uses fully supported coordinates;
different arbitrary extensions on unobserved directions
would not suffice.
Applying one common invertible coordinate change to both
contexts, including a common whitening, preserves the
contradiction.
Realigning the contexts separately changes the shared-map
question.

Testing a training-fitted family on held-out contexts
therefore asks more than whether each individual orbit
admits some action.
Success on finitely many held-out contexts still does not
establish a shared action on an unrestricted population of
contexts.


\subsection{Shared contextual realizations can still reflect state identity}

Shared realizability does not by itself identify semantic
computation.
Let $a_i\in\R^{1\times(n-1)}$ be a centered,
full-affine-rank state code, and let
$c_t\in\R^{1\times m}$ be a context vector.
By \cref{prop:finite}, there are matrices $B_g$ satisfying
$a_iB_g=a_{g(i)}$ and the row-action law.
For
\[
 z_{i,t}=(a_i,c_t),
 \qquad
 R_g=
 \begin{bmatrix}
  B_g&0\\
  0&\Id_m
 \end{bmatrix},
\]
one shared family maps all contexts exactly and obeys all
group relations.
The same state code supports alternative permutation
actions on the labels.
The construction requires no account of how a language
model computes a weekday shift or a geometric rotation.
A global invertible coordinate change preserves the
example up to conjugacy, so visible separation of state
and context coordinates is unnecessary.

The strict additive form has a further implication.
For a fixed invertible $T$, suppose
\[
 z_{i,t}=(a_i,c_t)T.
\]
Writing $\bar z_t=n^{-1}\sum_i z_{i,t}$ and
$\bar a=n^{-1}\sum_i a_i$, within-template centering gives
\[
 z_{i,t}-\bar z_t=(a_i-\bar a,0)T.
\]
All templates then share the same centered state code,
and the pooled centered matrix has rank at most $n-1$.
A fixed linear projection cannot increase that rank.

In \cref{app:template-centering-sensitivity}, all saved
centered Study~2 banks have numerical rank 16 under the
declared rule, exceeding the bounds 6 and 3 for this exact
weekday and corner construction.
The strict additive form therefore does not give an exact
description of those saved coordinates under that rank
criterion.
This observation neither measures its quality as an
approximation nor excludes stable state-code components
coexisting with context-dependent state variation.

Historical AFS-B changes operator labels inconsistently
with the prescribed multiplication table, while AFS-C
changes correspondences across training contexts;
both can damage a stable-code predictor.
The later category predictors and coherent alternative
registration in
\cref{sec:characterization,app:reproduction}
test different competing explanations.
Their predictive performance and the spectral comparisons
in \cref{sec:registration-sensitivity} address distinct
questions.
Neither identifies the strict construction above with
the measured activations or constitutes a conditional
comparison jointly matched in prediction error, operator
norm, information and capacity.


\subsection{Shared contextual actions as a subspace test}
\label{app:contextual-realization}

Stack contextual rows into $Z\in\R^{N\times q}$, with
complete state blocks, and let $P_g$ permute states within
each template.
For equal row weights and any rank of $Z$,
\[
 \min_{R\in\R^{q\times q}}
 \norm{ZR-P_gZ}_F^2
 =
 \norm{(\Id_N-ZZ^\dagger)P_gZ}_F^2.
\]
Indeed, $ZZ^\dagger P_gZ$ is the orthogonal projection of
the target matrix onto $\operatorname{col}(Z)$.
An exact shared map exists if and only if the residual
vanishes, equivalently
\[
 P_g\operatorname{col}(Z)
 \subseteq\operatorname{col}(Z).
\]
When $Z$ has full column rank, the exact solution
$R_g=Z^\dagger P_gZ$ is unique.
Existence for every group element then forces the
row-action multiplication law.

Equivalently, let
\[
 \mathcal N=\ker(Z^\top)
\]
be the space of linear dependencies among rows.
Exact realizability is equivalent to
\[
 P_g^\top\mathcal N\subseteq\mathcal N.
\]
Necessity follows from
\[
 a^\top Z=0
 \quad\Longrightarrow\quad
 a^\top P_gZ=a^\top ZR_g=0.
\]
Conversely, this condition makes
$a^\top Z\mapsto a^\top P_gZ$ well defined on the row span.
For a whole group, the induced maps are invertible and
obey the row-action law.

This is the finite row-observation form of the
kernel-preservation criterion for a transformation to
descend through a linear encoder
\citep[Theorem~1]{sun2026transformationlaws}.
The correspondence is explicit: take the linear map
$W=Z^\top$ and the input transformation $P_g^\top$, so that
\[
 R_g^\top W=WP_g^\top.
\]
Here $W$ acts on coefficients of observed rows; it is not
the language model's nonlinear encoder.
The condition concerns the left kernel of $Z$.
Preserving only equal-row pairs, or considering the right
kernel $\ker Z$, does not give this criterion.
For example, the three distinct scalar states
$(-1,0,1)$ have no exact linear realization of their
three-cycle, because $(0,1,-1)$ is not a scalar multiple
of $(-1,0,1)$.

When $Z$ is rank deficient, only the action on its row
span is identified.
Set
\[
 H_Z=Z^\dagger Z,
 \qquad
 S_g=Z^\dagger P_gZ.
\]
The minimum-norm identity solution is $S_e=H_Z$, a
projector rather than the ambient identity.
If all the subspace conditions hold, then
\[
 H_ZS_g=S_gH_Z=S_g
\]
and
\[
 \begin{aligned}
 S_gS_h
 &=Z^\dagger P_g(ZZ^\dagger)P_hZ\\
 &=Z^\dagger P_gP_hZ
 =S_{h\circ g}.
 \end{aligned}
\]
The second equality uses
$ZZ^\dagger P_hZ=P_hZ$.
Consequently,
\[
 \widetilde R_g=S_g+(\Id_q-H_Z)
\]
acts as the induced action on the row span and as identity
on its orthogonal complement, with
\[
 \widetilde R_e=\Id_q,
 \qquad
 \widetilde R_g\widetilde R_h
 =\widetilde R_{h\circ g}.
\]
Inverses follow from the group law.
The complementary action is a choice.
The action on the observed row span is unique; the full
ambient matrix is unique only when $Z$ has full column
rank.
New contexts can impose additional constraints.

For a complete uniformly weighted bank with full column
rank, exact linear realizability also connects directly
to the whitened class in \cref{eq:class-loss-main}.
Let
\[
 \Sigma=\frac1N Z^\top Z\succ0.
\]
If $ZR_g=P_gZ$, then
\[
 R_g^\top\Sigma R_g
 =\frac1N(P_gZ)^\top(P_gZ)
 =\Sigma.
\]
Thus
\[
 \rho_g=\Sigma^{1/2}R_g\Sigma^{-1/2}
\]
is orthogonal, realizes the action on $Z\Sigma^{-1/2}$,
and preserves the row-action law.
Under this common whitened loss,
\[
 \eta=0
 \quad\Longrightarrow\quad
 L_{\rm orth}=L^\star=0.
\]
This is a zero-risk implication.
It does not identify the three optima when their risks
are positive.
A weighted version requires the weighting matrix $D$ to
satisfy $P_g^\top DP_g=D$ and
$Z^\top DZ\succ0$.
Weights that vary only across complete templates satisfy
the invariance requirement.

The projection formula above is an unregularized,
unnormalized, finite-bank characterization.
It is not the row-normalized held-out score in
\cref{eq:D}, and it gives no population guarantee.
Orthogonality in the original Euclidean coordinates or
in another specified metric requires its own check.
A rank-deficient bank also does not satisfy the
$q$-dimensional unit-second-moment condition merely
because an invertible ambient extension exists.

For comparison, the ridge objective
\[
 \norm{ZR-P_gZ}_F^2+\lambda\norm R_F^2,
 \qquad \lambda>0,
\]
returns
\[
 R_g^{(\lambda)}
 =(Z^\top Z+\lambda\Id_q)^{-1}Z^\top P_gZ.
\]
Here $\lambda$ is the penalty in this unnormalized
objective; a coefficient specified in an
energy-normalized implementation must first be converted
to this convention.
If an exact $R_g^\star$ exists, then
\[
 R_g^{(\lambda)}=B_\lambda R_g^\star,
 \qquad
 B_\lambda
 =(Z^\top Z+\lambda\Id_q)^{-1}Z^\top Z.
\]
In particular, the independently fitted identity is
$B_\lambda$, which generally differs from $\Id_q$.
Shrinkage can introduce power and relation errors even
in an exactly realizable system.
The relation and identity penalties in Joint address
these departures, but their empirical effects do not
by themselves isolate additional semantic information.


\subsection{An exact common-design ridge identity}
\label{app:ridge-identity}

Suppose every element uses the same source matrix
$Z\in\R^{N\times q}$, the same effective penalty
$\lambda>0$, and targets $P_gZ$ from one coherent group
action.
Define
\begin{align*}
 L_\lambda
 &=(Z^\top Z+\lambda\Id_q)^{-1}Z^\top,\\
 \mathsf E_\lambda
 &=\Id_N-ZL_\lambda,
 \qquad
 R_g=L_\lambda P_gZ.
\end{align*}
The symbol $\mathsf E_\lambda$ denotes the ridge residual
operator, not an orbit-Gram matrix.

Using $P_gP_h=P_{h\circ g}$ gives
\begin{align}
 ZR_g-P_gZ
 &=-\mathsf E_\lambda P_gZ,
 \label{eq:ridge-fit-identity}\\
 R_gR_h-R_{h\circ g}
 &=-L_\lambda P_g\mathsf E_\lambda P_hZ.
 \label{eq:ridge-law-identity}
\end{align}
For the first identity, substitute
$ZL_\lambda=\Id_N-\mathsf E_\lambda$ into $ZR_g$.
For the second,
\[
 \begin{aligned}
 R_gR_h
 &=L_\lambda P_g(ZL_\lambda)P_hZ\\
 &=L_\lambda P_gP_hZ
   -L_\lambda P_g\mathsf E_\lambda P_hZ.
 \end{aligned}
\]
No exact realizability or full-rank assumption is needed.
For $\lambda>0$, $\mathsf E_\lambda$ is generally not
an orthogonal projector.

Coherence of the target action is essential to the
second identity.
If instead the targets use arbitrary permutation matrices
$\widetilde P_g$, with
$\widetilde R_g=L_\lambda\widetilde P_gZ$, then
\[
 \begin{aligned}
 \widetilde R_g\widetilde R_h
 -\widetilde R_{h\circ g}
 &=-L_\lambda\widetilde P_g
      \mathsf E_\lambda\widetilde P_hZ\\
 &\quad+
 L_\lambda
 \bigl(\widetilde P_g\widetilde P_h
       -\widetilde P_{h\circ g}\bigr)Z.
 \end{aligned}
\]
The extra term prevents applying
\cref{eq:ridge-law-identity} unchanged to inconsistent
operator relabelings or target scrambling.

The common-design identities apply directly to the full
independent Ridge fit in Study~1.
They do not apply unchanged to Joint after optimization,
to a family formed by composing fitted generators, or
to fits with element-specific source subsets, weights,
penalties or independently centered targets.
They explain an algebraic dependence between fit and
composition defects; they are not independent
significance tests.
Their unnormalized matrix residuals must also be
distinguished from the row-normalized test metric in
\cref{eq:D}.


\subsection{Equal within-context geometry need not permit shared transport}
\label{app:orientation}

We give a restricted analytical construction using
orthogonal representations and group averaging
\citep{serre1977linear}.
It specifies an identification boundary, rather than
a model assumed for the measured activations.

Let $1\leq k\leq n-1$ and let
$V\in\R^{n\times k}$ satisfy
\[
 \mathbf1^\top V=0,
 \qquad
 V^\top V/n=\Id_k,
 \qquad
 VQ_g=P_gV,
\]
where $Q_g\in O(k)$ and
$Q_gQ_h=Q_{h\circ g}$.
For $d\geq k$, choose
$A_c\in\R^{d\times k}$ with
$A_c^\top A_c=\Id_k$, and set
\[
 X_c=VA_c^\top.
\]
The frames belong to the Stiefel manifold
\[
 \mathrm{St}(d,k)
 =\{A\in\R^{d\times k}:A^\top A=\Id_k\}
\]
\citep{edelman1998geometry}.
This is a space of frame matrices in the construction,
not an assumed manifold of language-model activations.
Here $k$ is the intrinsic orbit rank and $d$ is the
ambient observation dimension.

Every context has the same within-context Gram matrix,
\[
 X_cX_c^\top=VV^\top,
\]
and an exact local orthogonal action
\[
 R_{c,g}
 =A_cQ_gA_c^\top+(\Id_d-A_cA_c^\top).
\]
Indeed,
\[
 X_cR_{c,g}=P_gX_c,
\]
and the two orthogonal subspaces carry $Q_g$ and the
identity, respectively.
The construction does not fix cross-context inner
products.

\begin{proposition}[Independent-orientation risk]
\label{prop:orientation}
For a nonempty $\mathcal A\subseteq G$, define
\[
 \ell_c(R)
 =\frac{1}{nk|\mathcal A|}
   \sum_{g\in\mathcal A}
   \norm{X_cR_g-P_gX_c}_F^2.
\]
Suppose a fresh $A_c$ has the uniform Stiefel distribution
induced by Haar measure on $O(d)$ and is independent of
the fitted matrices $\widehat R_g\in\R^{d\times d}$.
Then
\begin{equation}
\begin{aligned}
 &\mathbb E[\ell_c(\widehat R)\mid\widehat R]\\
 &\quad=
 b_{\mathcal A}
 +\frac{1}{d|\mathcal A|}
  \sum_{g\in\mathcal A}
  \norm{\widehat R_g-\frac{\chi(g)}k\Id_d}_F^2,
\end{aligned}
\end{equation}
where
\[
 \chi(g)=\operatorname{tr}(Q_g),
 \qquad
 b_{\mathcal A}
 =1-\frac1{|\mathcal A|}
       \sum_{g\in\mathcal A}
       \left(\frac{\chi(g)}k\right)^2.
\]
Thus $b_{\mathcal A}$ is the minimum expected risk among
unrestricted shared matrix predictors in this
construction.
\end{proposition}

\begin{proof}
Using $V^\top V=n\Id_k$,
\[
 \frac{\norm{X_cR-P_gX_c}_F^2}{nk}
 =
 \frac{\norm{A_c^\top R-Q_gA_c^\top}_F^2}{k}.
\]
Orthogonal invariance and traces give
\[
 \mathbb E[AA^\top]=\frac{k}{d}\Id_d,
 \qquad
 \mathbb E[AQ_gA^\top]=\frac{\chi(g)}d\Id_d.
\]
Expanding the square and taking expectations yields
\begin{align*}
 &\frac1d\norm R_F^2+1
   -\frac{2\chi(g)}{kd}\operatorname{tr}R\\
 &\quad=
 1-\frac{\chi(g)^2}{k^2}
 +\frac1d
  \norm{R-\frac{\chi(g)}k\Id_d}_F^2.
\end{align*}
Averaging over $\mathcal A$ proves the identity.
Independence permits conditioning on the fitted family.
The scalar matrices
$R_g=\chi(g)\Id_d/k$ attain the minimum, although they
need not form a group representation.
\end{proof}

The baseline $b_{\mathcal A}$ is an unconstrained
population prediction risk.
It is distinct from the orthogonality cost
$\beta=L_{\rm orth}-\eta$ in the main text.

For the full centered permutation representation of the
seven-state cycle,
\[
 k=6,\qquad
 \chi(e)=6,\qquad
 \chi(g)=-1\quad(g\ne e).
\]
Hence
\[
 b_G=\frac56,
 \qquad
 b_{G\setminus\{e\}}=\frac{35}{36}.
\]
For the full centered four-corner permutation
representation of $D_4$, $k=3$, and every nonidentity
character is $+1$ or $-1$, giving
\[
 b_G=\frac79,
 \qquad
 b_{G\setminus\{e\}}=\frac89.
\]
These calculations use the complete centered permutation
representations, not two-dimensional polygon coordinates.

The normalization above uses the orbit energy $nk$.
For one context,
\[
 \frac1nX_c^\top X_c=A_cA_c^\top,
\]
which has rank $k$ and is singular when $d>k$.
Under the specified Haar distribution,
\[
 \mathbb E_c\!\left[\frac1nX_c^\top X_c\right]
 =\frac{k}{d}\Id_d,
\]
but a finite context bank need not have this second
moment or even full ambient rank.
The constants therefore do not directly provide bounds
for an empirical bank whitened by its own second moment.

The full orthogonal-invariance assumption is essential.
Independence alone, or a distribution restricted to a
commuting family of rotations, is insufficient.
The constants also cannot be compared numerically with
the reported $\Eeq$, which uses a different distribution
and normalization.
No Haar randomization test or empirical rejection of
this construction is claimed.


\subsection{A commutant criterion in the full-rank construction}
\label{app:commutant}

Specialize the preceding construction to $d=k$, writing
$A_c=O_c\in O(k)$.
The orientations may now have any distribution.
Define
\begin{align*}
 \mathcal C_Q
 &=\{B\in\R^{k\times k}:BQ_g=Q_gB
                         \text{ for every }g\in G\},\\
 \Pi_Q(B)
 &=\frac1{|G|}\sum_{g\in G}Q_gBQ_g^\top.
\end{align*}
The group average is the Frobenius orthogonal projection
onto $\mathcal C_Q$.
Reindexing shows that its image commutes with every
$Q_g$; it fixes every commuting matrix; and inverse
closure makes it self-adjoint.

\begin{proposition}[Relative-orientation obstruction]
\label{prop:commutant}
Let $c,c'$ be independent draws from the context
distribution, set $U=O_c^\top O_{c'}$, and use
$\mathcal A=G$ in $\ell_c$.
Then
\[
 \inf_{\substack{R_g\in\R^{k\times k}\\g\in G}}
 \mathbb E_c\ell_c(R)
 =
 \frac1k
 \mathbb E_{c,c'}
 \norm{U-\Pi_Q(U)}_F^2.
\]
An exact shared action exists almost surely if and only
if the relative orientations commute with every $Q_g$
almost surely.
\end{proposition}

\begin{proof}
The local action is
\[
 R_{c,g}=O_cQ_gO_c^\top.
\]
Since $X_c^\top X_c/n=\Id_k$, the loss for element $g$
equals
\[
 \frac1k\norm{R_g-R_{c,g}}_F^2.
\]
Its unrestricted minimizer is
$\mathbb E_cR_{c,g}$.
The pairwise variance identity therefore gives the
optimum risk as
\[
 \frac1{2k|G|}
 \mathbb E_{c,c'}
 \sum_{g\in G}\norm{Q_gU-UQ_g}_F^2.
\]
Expansion and the projection identity yield
\[
 \frac1{|G|}
 \sum_{g\in G}\norm{Q_gU-UQ_g}_F^2
 =
 2\norm{U-\Pi_Q(U)}_F^2.
\]
This proves the risk formula.
Nonnegativity gives the zero-risk equivalence.
At zero risk, the common matrices agree almost surely
with the local orthogonal actions and inherit their
group law.
At positive risk, their means need not be orthogonal
or form a representation.
\end{proof}

Large context rotations within the commutant can leave
the action unchanged.
The proposition is a characterization of this
full-rank, whitened, noiseless construction.
Its infimum is over unrestricted shared matrices,
rather than the exact-action class defining $L^\star$.

There is also a deterministic finite-bank version.
For any fixed $m\geq1$ contexts, define
\[
 \widehat R_{m,g}
 =\frac1m\sum_{j=1}^mR_{c_j,g},
 \qquad
 \widehat L_m
 =\frac1m\sum_{j=1}^m\ell_{c_j}(\widehat R_m).
\]
The sample mean minimizes the unrestricted bank loss.
Writing $U_{ij}=O_{c_i}^\top O_{c_j}$, the same pairwise
variance identity gives
\[
 \widehat L_m
 =
 \frac1{km^2}
 \sum_{i=1}^m\sum_{j=1}^m
 \norm{U_{ij}-\Pi_Q(U_{ij})}_F^2.
\]
The sums include all ordered pairs.
No independence assumption is needed for this fixed-bank
identity.
Because every context has unit second moment, this is
the independent-linear optimum $\eta$ for the same
uniformly weighted bank and full-group loss.
It does not determine $L^\star$ when the risk is positive,
and zero bank risk alone gives no guarantee for new
contexts.

For iid training contexts, let
\[
 b=
 \inf_{\substack{R_g\in\R^{k\times k}\\g\in G}}
 \mathbb E_c\ell_c(R).
\]
If $c_{\rm new}$ is an independent draw from the same
distribution, the variance identity gives
\[
 \mathbb E\widehat L_m=(1-1/m)b,
 \qquad
 \mathbb E\ell_{c_{\rm new}}(\widehat R_m)
 =(1+1/m)b.
\]
Both expectations include the sampled training contexts;
the second also includes the fresh context.
These sample-mean identities are analytical consequences
of this construction, not learning-curve predictions
for ridge fitted to the actual 16-dimensional activations.


\subsection{Prediction-error propagation}
\label{app:propagation-proof}

\begin{proof}[Proof of \cref{prop:propagation}]
Insert and subtract
$z_{g_2(i),t}R_{g_1}$ and
$z_{g_1(g_2(i)),t}$.
Since $p(i)=g_1(g_2(i))$,
\begin{align*}
 z_{i,t}R_{g_2}R_{g_1}-z_{i,t}R_p
 &=\bigl(z_{i,t}R_{g_2}-z_{g_2(i),t}\bigr)R_{g_1}\\
 &\quad+
 \bigl(z_{g_2(i),t}R_{g_1}-z_{p(i),t}\bigr)\\
 &\quad-
 \bigl(z_{i,t}R_p-z_{p(i),t}\bigr).
\end{align*}
This is the stated identity.
The triangle inequality and
$\norm{vA}_2\leq\norm v_2\norm A_{\rm op}$
give the norm bound.
\end{proof}

For example, suppose every unnormalized prediction error
is at most $\delta$ on complete, group-closed evaluation
state blocks, and
$\norm{R_g}_{\rm op}\leq M$ for every relevant element.
Then every pairwise residual is at most
$(M+2)\delta$.
If the comparison denominators also satisfy
\[
 \norm{z_{i,t}R_p}_2^2+10^{-12}
 \geq b_{\min}>0,
\]
each corresponding normalized defect is at most
\[
 \frac{(M+2)^2\delta^2}{b_{\min}}.
\]

The uniform error assumption must cover all elements
used in the identity, including the fitted $R_e$ when
$p=e$.
The prediction aggregate $\Eeq$ excludes the identity,
whereas the products in the consistency diagnostics can
equal it.
Moreover, the consistency denominator comes from
$z_{i,t}R_p$, the directly fitted prediction, rather
than the true target $z_{p(i),t}$.

A small average prediction error supplies neither a
uniform $\delta$ nor a useful denominator lower bound.
Changed-state-only aggregates likewise do not provide
uniform control over a complete group-closed set.
The identity therefore does not yield a numerical bound
on the reported $D$ without accounting for its
sample-specific denominators and operator norms.
It explains a dependence that matters when assessing
algebraic consistency beyond predictive accuracy.
The present experiments do not include an error-matched
conditional null analysis.


\subsection{Global additive cancellation}
\label{app:global-cancellation-proof}

\begin{proof}[Proof of \cref{prop:additive}]
Since $g$ is a bijection of $S$,
\[
 \sum_{i\in S}z_{g(i),t}
 =
 \sum_{i\in S}z_{i,t}.
\]
Subtracting gives zero separately in every template.
Averaging over any common distribution of training
templates preserves zero.
\end{proof}

The cancellation holds in every representation dimension
and does not require centering or PCA.
It concerns the sum of transformation displacements:
\[
 \sum_i(z_{g(i),t}-z_{i,t})=0.
\]
It does not assert that the template mean
$n^{-1}\sum_i z_{i,t}$ is zero.
Thus it does not imply removal of the state-invariant
component studied in
\cref{app:template-centering-sensitivity}.
A numerical implementation can retain roundoff residuals.

More generally, cancellation holds for weights invariant
under the relevant permutation.
For a fixed $g$, deleting fixed-point pairs also preserves
cancellation whenever moved states remain, because each
deleted displacement is zero.
Incomplete state sets that are not preserved by $g$, or
non-invariant weights, need not cancel.
Those conditions differ from the complete balanced
training estimator considered here.

A source-state-conditioned vector,
\[
 v_{g,i}
 =\mathbb E_t[h_{g(i),t}-h_{i,t}],
\]
need not vanish because it does not average over source
states.


\subsection{State-conditioned additive actions can compose}
\label{app:additive-composition}

For training-only state means
\[
 \mu_i
 =\mathbb E_{t\in\mathcal T_{\mathrm{train}}}h_{i,t},
\]
using the same template weights for every state, the
source-conditioned comparator satisfies
\[
 v_{g,i}
 =\mathbb E_t(h_{g(i),t}-h_{i,t})
 =\mu_{g(i)}-\mu_i.
\]
To compose the edits, retain the state index and update
it after each action.
Define the augmented-state function
\[
 \mathcal A_g(h,i)
 =
 \bigl(h+\mu_{g(i)}-\mu_i,\ g(i)\bigr).
\]
For $p=g_1\circ g_2$,
\begin{align*}
 (\mathcal A_{g_1}\circ\mathcal A_{g_2})(h,i)
 &=
 \bigl(
 h+(\mu_{g_2(i)}-\mu_i)
   +(\mu_{p(i)}-\mu_{g_2(i)}),
 p(i)
 \bigr)\\
 &=
 \bigl(h+\mu_{p(i)}-\mu_i,p(i)\bigr)\\
 &=\mathcal A_p(h,i).
\end{align*}
The means telescope, and
$\mathcal A_e(h,i)=(h,i)$.
This is an exact action on the augmented space.
For noncommuting actions, the two orders can have
different target means, so the construction can preserve
order as well as composition.

The state index is part of the construction.
Without that index, this argument does not define a
linear operator family on the original hidden space.
The historical behavioral experiment supplies the source
identity as oracle information and fits displacement
vectors in the full hidden dimension.
It therefore differs from a structural comparison in
the same PCA coordinates with matched identity
information.
The coordinate-level category controls in
\cref{sec:characterization,app:reproduction} retain their
own stated information conditions.
None of these constructions identifies a native neural
mechanism.

The exact composition guarantee also differs from
re-decoding every intermediate vector.
For nearest-centroid decoding $\widehat i(h)$, define
the stateless rule
\[
 F_g(h)
 =
 h+\mu_{g(\widehat i(h))}-\mu_{\widehat i(h)}.
\]
With $\mu_0=0$, $\mu_1=2$ and the swap $s$,
\[
 F_s(-2)=0,\qquad F_s(0)=2.
\]
Thus
\[
 F_s(F_s(-2))=2\ne-2,
\]
although $s^2=e$.
Perfect decoding on clean evaluation inputs does not
certify decoding after edits.
Likewise, looking up a known composite action after one
input decode does not test repeated learned primitive
edits.

Global additive cancellation consequently does not
establish an intrinsic advantage of multiplicative
notation.
Every linear predictor can also be written as a
state-dependent displacement:
\[
 zR_g=z+z(R_g-\Id).
\]


\subsection{Faithful actions need not be free}
\label{app:faithful-not-free}

An action is faithful if no nonidentity group element
fixes every state.
It is free if no nonidentity element fixes any state.
The four-corner $D_4$ action is faithful but not free:
each corner has a stabilizer of size two.
A fixed source corner therefore has four possible
endpoints for eight group elements.

The complete labeled state correspondence distinguishes
the elements of this faithful action.
An individual source--target pair does not.
This statement concerns state labels; hidden
representations that merge states need not preserve
that distinction.
Faithfulness of the prescribed input action also does
not restrict the class in \cref{eq:class-loss-main} to
faithful representation maps.

Number the corners clockwise as
\[
 0=\mathrm{NW},\qquad
 1=\mathrm{NE},\qquad
 2=\mathrm{SE},\qquad
 3=\mathrm{SW}.
\]
Then the generators in \cref{sec:setup} are
\[
 r(i)=i+1,\qquad s(i)=1-i
 \pmod 4.
\]
Words here denote composition of input functions, with
the rightmost function applied first.
Their corresponding row-matrix products have the
reversed order.
We have
\[
 srs(i)=i-1,\qquad rsr(i)=1-i,
\]
which agree at $i=1,3$, whereas
\[
 rs(i)=2-i,\qquad sr(i)=-i
\]
differ at every corner.
Order comparisons must therefore report which source
states distinguish the tested endpoints.
Success on this four-state action does not establish
a free eight-state action or native group multiplication.


\subsection{Coordinate equivalence and the metric of approximation}
\label{app:coordinates}

For an invertible coordinate change $T_{\mathrm{coord}}$,
set
\[
 z'=zT_{\mathrm{coord}},
 \qquad
 R'_g=T_{\mathrm{coord}}^{-1}R_gT_{\mathrm{coord}}.
\]
For $p=g_1\circ g_2$,
\[
 \begin{aligned}
 z'R'_g-z'_g
 &=(zR_g-z_g)T_{\mathrm{coord}},\\
 R'_{g_2}R'_{g_1}-R'_p
 &=
 T_{\mathrm{coord}}^{-1}
 (R_{g_2}R_{g_1}-R_p)
 T_{\mathrm{coord}}.
 \end{aligned}
\]
Exact realization and exact group identities are
therefore invariant under simultaneous conjugacy.
An orthogonal coordinate change also preserves the
Euclidean row norms and Frobenius norms used in the
recorded defects.
With Euclidean scoring in both coordinate systems, a
general invertible change can alter approximate errors
and the weighting of directions.
Euclidean ridge penalties are likewise coordinate
dependent, so refitting need not return the conjugated
original estimate.

A different conclusion applies when the same full-rank
bank is reparameterized and then whitened by its own
second moment.
Let
\[
 \Sigma=\mathbb E[z^\top z]\succ0,
 \qquad
 z'=zT,
 \qquad
 \Sigma'=T^\top\Sigma T.
\]
Define
\[
 x=z\Sigma^{-1/2},
 \qquad
 x'=zT(\Sigma')^{-1/2},
 \qquad
 O_T=\Sigma^{1/2}T(\Sigma')^{-1/2}.
\]
Then
\[
 O_T^\top O_T=\Id,
 \qquad
 x'=xO_T.
\]
The two whitened banks differ only by an orthogonal
coordinate change.
Conjugating a family by $O_T$ preserves its loss,
orthogonality and group law.
Consequently, with the same observations, state
correspondences and weights, the corresponding
$\eta$, $L_{\rm orth}$, $L^\star$ and type-constrained
spectral budgets are unchanged.
Candidate-image rankings are unchanged as well.

This invariance does not identify the three treatments
in \cref{app:template-centering-sensitivity}.
Within-template centering changes the observations,
and changing the retained PCA subspace is not an
invertible change of coordinates on the same space.
Moreover, centered observations retaining the original
whitening generally do not have unit second moment.
Their score $C_{16}/16$ uses the original energy scale;
the invariance between two self-whitened descriptions
does not equate that score with the rewhitened one.

The prescribed group action on input states is also
distinct from a coordinate-equivalence group acting on
the representation description
\citep{yao2026groupaction}.
Context-specific realignments change the shared-map
question.
Changing the state correspondence while keeping the
representations fixed defines another registered task.
These operations must be distinguished when interpreting
semantic preference.

For an unnormalized activation-conditioned matrix
defect $F$,
\[
 \mathbb E\norm{zF}_2^2
 =
 \operatorname{tr}(F^\top\Sigma F),
 \qquad
 \Sigma=\mathbb E[z^\top z].
\]
Here $\Sigma$ is the uncentered second moment.
The identity shows how observation energy weights
directions.
A small data-space defect can reflect low energy in a
direction even when that direction is not exactly
unsupported.
The actual $D(A,B)$ additionally normalizes each row by
its comparison-target energy.
The operator diagnostics use different powers and, for
the C7 relation diagnostic, different aggregands.
Their numerical comparison therefore does not by itself
establish a sharp supported/unsupported distinction.

For any exact finite group action by invertible matrices,
a positive-definite invariant row metric exists.
Specifically,
\[
 H=\frac1{|G|}\sum_{g\in G}R_gR_g^\top
\]
is positive definite and satisfies
\[
 R_hHR_h^\top=H.
\]
The equality follows because multiplication by $R_h$
permutes the group-indexed family.
Choose $T_{\mathrm{coord}}$ with
$T_{\mathrm{coord}}T_{\mathrm{coord}}^\top=H$.
Then
\[
 \begin{aligned}
 R'_g(R'_g)^\top
 &=
 T_{\mathrm{coord}}^{-1}
 R_gHR_g^\top
 T_{\mathrm{coord}}^{-\top}\\
 &=\Id.
 \end{aligned}
\]
This construction chooses a metric for the given exact
action.
It does not make the original Euclidean metric invariant
or preserve a comparison that has already fixed a
different metric.
The complete-bank zero-risk argument in
\cref{app:contextual-realization} separately explains
when the bank's own whitening supplies an orthogonal
realization.
Approximately consistent fitted matrices need not
satisfy the invariant-metric identity above.

Finally, character references must match the represented
space.
The $(n-1)$-dimensional centered permutation
representation of \cref{prop:finite} has character
\[
 \chi(g)=\#\{i\in S:g(i)=i\}-1,
\]
whereas the full permutation representation has
character equal to the number of fixed states.
Additional state-invariant coordinates contribute
additional trivial characters.
Centering by the overall training mean does not
automatically remove the state-invariant variation
across templates.

These formulas therefore do not prescribe the character
of an arbitrary centered 16-dimensional representation.
Comparing normalized traces with a fixed-state-count
reference does not measure distance to every valid
representation.
Character-inner-product expressions computed from
approximately consistent fitted matrices remain
descriptive real numbers; they are not certified
integer irreducible multiplicities.
